% 48e Olympiade Internationale de Mathématiques — Hanoï, Viêt Nam.
% Énoncés officiels en français (imo-official.org). Chaque problème vaut 7 points.

\begin{center}
  \emph{Premier jour --- mercredi 25 juillet 2007 --- Durée : 4 heures 30}
\end{center}

\subsection*{Problème 1}
Soit $n$ nombres réels $a_1, a_2, \ldots, a_n$. Pour chaque $i$ ($1 \leq i \leq n$) on définit
\[
d_i = \max\{a_j : 1 \leq j \leq i\} - \min\{a_j : i \leq j \leq n\}
\]
et on pose
\[
d = \max\{d_i : 1 \leq i \leq n\}.
\]
\begin{enumerate}
  \item[(a)] Montrer que pour tous nombres réels $x_1 \leq x_2 \leq \cdots \leq x_n$,
  \[
  \max\{|x_i - a_i| : 1 \leq i \leq n\} \geq \frac{d}{2}. \tag{$*$}
  \]
  \item[(b)] Montrer qu'il existe des nombres réels $x_1 \leq x_2 \leq \cdots \leq x_n$ tels que $(*)$ soit
  une égalité.
\end{enumerate}

\subsection*{Problème 2}
On donne cinq points $A, B, C, D$ et $E$ tels que $ABCD$ soit un parallélogramme et $BCED$ un quadrilatère
convexe, inscriptible. Soit $\ell$ une droite passant par $A$. On suppose que $\ell$ coupe l'intérieur du
segment $DC$ en $F$ et coupe la droite $BC$ en $G$. On suppose aussi que $EF = EG = EC$. Montrer que $\ell$ est
la bissectrice de l'angle $\widehat{DAB}$.

\subsection*{Problème 3}
Dans une compétition mathématique certains participants sont des amis. L'amitié est toujours réciproque. Un
groupe de participants est appelé une \emph{clique} si toute paire d'entre eux est formée de deux amis. (En
particulier, chaque groupe d'au plus un participant constitue une clique.) Le nombre de participants dans une
clique est appelé sa \emph{taille}.

On suppose que, dans cette compétition, la plus grande taille des cliques est paire. Montrer que les
participants peuvent être répartis dans deux pièces de telle sorte que la plus grande taille des cliques
contenues dans une de ces pièces soit égale à la plus grande taille des cliques contenues dans l'autre.

\bigskip
\begin{center}
  \emph{Second jour --- jeudi 26 juillet 2007 --- Durée : 4 heures 30}
\end{center}

\subsection*{Problème 4}
Dans un triangle $ABC$ la bissectrice de l'angle $\widehat{BCA}$ recoupe le cercle circonscrit en $R$, coupe la
médiatrice de $BC$ en $P$ et la médiatrice de $AC$ en $Q$. Le milieu de $BC$ est $K$ et le milieu de $AC$ est
$L$. Montrer que les triangles $RPK$ et $RQL$ ont la même aire.

\subsection*{Problème 5}
Soit $a$ et $b$ deux entiers strictement positifs. Montrer que si $4ab - 1$ divise $(4a^2 - 1)^2$, alors
$a = b$.

\subsection*{Problème 6}
Soit $n$ un entier strictement positif. Dans l'espace on considère l'ensemble
\[
S = \bigl\{(x, y, z) \ : \ x, y, z \in \{0, 1, \ldots, n\},\ x + y + z > 0\bigr\},
\]
constitué de $(n + 1)^3 - 1$ points. Trouver le plus petit nombre de plans dont la réunion contient $S$ mais ne
contient pas $(0, 0, 0)$.
